\documentclass[12pt, a4paper]{article}

\usepackage{amsmath, amssymb, amsthm}
\usepackage{geometry}
\usepackage{fancyhdr}
\usepackage{enumitem}
\usepackage{xcolor}
\usepackage{hyperref}
\usepackage{tcolorbox}
\usepackage{xeCJK}
% ── 中文字体自动探测（由 render_latex.py 生成）──
% 探测到的候选字体：Source Han Serif SC, Noto Sans SC, SimSun, Microsoft YaHei, SimHei, KaiTi
\IfFontExistsTF{Source Han Serif SC}{\setCJKmainfont{Source Han Serif SC}}{}
\IfFontExistsTF{Noto Sans SC}{\setCJKmainfont{Noto Sans SC}}{}
\IfFontExistsTF{SimSun}{\setCJKmainfont{SimSun}}{}
\IfFontExistsTF{Microsoft YaHei}{\setCJKmainfont{Microsoft YaHei}}{}
\IfFontExistsTF{SimHei}{\setCJKmainfont{SimHei}}{}
\IfFontExistsTF{KaiTi}{\setCJKmainfont{KaiTi}}{}
\IfFontExistsTF{FangSong}{\setCJKmainfont{FangSong}}{}
\IfFontExistsTF{STSong}{\setCJKmainfont{STSong}}{}
% 若以上全部缺失，中文可能显示为方框；请安装 Noto Sans CJK SC 或在 config 中指定字体
\XeTeXlinebreaklocale "zh"
\XeTeXlinebreakskip = 0pt plus 1pt
% 关闭 CJK 与西文之间的额外间距，贴近学术排版习惯

\geometry{margin=1in}
\pagestyle{fancy}
\fancyhf{}
\lhead{Mathematics}
\rhead{Student}
\rfoot{Page \thepage}
\cfoot{}

% ── 颜色定义 ──
\definecolor{solutionblue}{HTML}{1D4ED8}
\definecolor{solutiongreen}{HTML}{15803D}
\definecolor{warnamber}{HTML}{B45309}
\definecolor{verifyred}{HTML}{B91C1C}

% ── 自定义环境 ──
\newtcolorbox{stepbox}{colback=blue!2!white,colframe=blue!20!black,boxrule=0.6pt,arc=2pt,left=6pt,right=6pt,top=4pt,bottom=4pt}
\newtcolorbox{verifybox}{colback=green!3!white,colframe=green!35!black,boxrule=0.6pt,arc=2pt,left=6pt,right=6pt,top=3pt,bottom=3pt}
\newtcolorbox{warnbox}{colback=orange!5!white,colframe=orange!45!black,boxrule=0.6pt,arc=2pt,left=6pt,right=6pt,top=3pt,bottom=3pt}

\newcommand{\problem}[1]{\subsection*{Problem~#1}}
\newcommand{\solution}{%
  \par\noindent
  {\color{solutiongreen}\bfseries Solution.}\par
}
\newcommand{\verifyok}{\textcolor{solutiongreen}{\ensuremath{\checkmark}\ 已验证}}
\newcommand{\verifywarn}{\textcolor{verifyred}{\ensuremath{\times}\ 验证存疑}}

\begin{document}

\begin{center}
  {\LARGE\bfseries Homework Solutions} \\[0.6em]
  {\large\color{solutionblue} Mathematics} \\[0.4em]
  Student \quad | \quad 2026-10-06
\end{center}
\vspace{0.5em}\hrule\vspace{1.2em}
\problem{Q1}

\begin{stepbox}
Can the graph of a function have more than one tangent at a given point? If so, graph your answer. If not, explain why.
\end{stepbox}

\solution

\begin{stepbox}
**不能**——在可导点处切线**唯一**。
\end{stepbox}

\begin{stepbox}
理由：切线斜率由导数唯一确定，$k=f'(a)$；
\end{stepbox}

\begin{stepbox}
过点 $(a,f(a))$ 且斜率为 $k$ 的直线只有一条。
\end{stepbox}

\begin{stepbox}
（若导数不存在，则是「零条」切线而非「多条」。）
\end{stepbox}

\begin{center}
\textbackslash{}text\{唯一切线\}:\textbackslash{} y-f(a)=f'(a)(x-a)
\end{center}

\begin{verifybox}\small \verifyok\quad 1/1 项检验通过（答案形式检查） \end{verifybox}

\vspace{0.8em}\hrule\vspace{0.8em}

\problem{Q2}

\begin{stepbox}
Is there a function whose graph doesn't have a tangent at some point? If so, graph your answer. If not, explain why.
\end{stepbox}

\solution

\begin{stepbox}
**是**，存在函数在某点没有切线。
\end{stepbox}

\begin{stepbox}
反例：$f(x)=|x|$ 在 $x=0$ 处连续但不可导——
\end{stepbox}

\begin{stepbox}
左右导数分别为 $f'_-(0) = -1$ 与 $f'_+(0) = +1$，不相等。
\end{stepbox}

\begin{stepbox}
图像在原点处有一个「尖点」，不存在切线。
\end{stepbox}

\begin{center}
f(x)=|x| \textbackslash{}text\{ 在 \} x=0 \textbackslash{}text\{ 处不可导\}\textbackslash{} (f'\_-(0)=-1\textbackslash{}neq f'\_+(0)=1)
\end{center}

\begin{verifybox}\small \verifyok\quad 1/1 项检验通过（答案形式检查） \end{verifybox}

\vspace{0.8em}\hrule\vspace{0.8em}

\problem{P1}

\begin{stepbox}
Suppose $\delta$ is a positive real number. How do you describe all real numbers $x$ that are within $\delta$ of 0?
\end{stepbox}

\solution

\begin{stepbox}
x 在 0 的 δ 邻域内：$|x - 0| < \delta$，即 $0-\delta < x < 0+\delta$
\end{stepbox}

\begin{center}
$\boxed{|x - 0| < \delta}$
\end{center}

\begin{verifybox}\small \verifyok\quad 1/1 项检验通过（答案形式检查） \end{verifybox}

\vspace{0.3em}
\textbf{(a)}\quad \textcolor{gray}{\small Using inequalities (\$<, >\$).}

用不等式表示：$0 - \delta < x < 0 + \delta$

$\quad\Rightarrow\quad 0 - \delta < x < 0 + \delta$

\textbf{(b)}\quad \textcolor{gray}{\small Using absolute value notation.}

$|x - 0| < \delta$

$\quad\Rightarrow\quad |x - 0| < \delta$

\textbf{(c)}\quad \textcolor{gray}{\small Using interval notation.}

区间表示：$(0 - \delta, 0 + \delta)$

$\quad\Rightarrow\quad (0 - \delta, 0 + \delta)$


\vspace{0.8em}\hrule\vspace{0.8em}

\problem{P2}

\begin{stepbox}
Suppose $a$ is a real number and $\delta$ is a positive real number. How do you describe all real numbers $x$ that are within $\delta$ of $a$?
\end{stepbox}

\solution

\begin{stepbox}
x 在 0 的 δ 邻域内：$|x - 0| < \delta$，即 $0-\delta < x < 0+\delta$
\end{stepbox}

\begin{center}
$\boxed{|x - 0| < \delta}$
\end{center}

\begin{verifybox}\small \verifyok\quad 1/1 项检验通过（答案形式检查） \end{verifybox}

\vspace{0.3em}
\textbf{(a)}\quad \textcolor{gray}{\small Graphically, as on the line above.}

\textcolor{warnamber}{未解：未识别的 ε-δ 子题类型}

\textbf{(b)}\quad \textcolor{gray}{\small Using inequalities. There are at least two ways to do this. Find both.}

用不等式表示：$0 - \delta < x < 0 + \delta$

$\quad\Rightarrow\quad 0 - \delta < x < 0 + \delta$

\textbf{(c)}\quad \textcolor{gray}{\small Using absolute value notation.}

$|x - 0| < \delta$

$\quad\Rightarrow\quad |x - 0| < \delta$

\textbf{(d)}\quad \textcolor{gray}{\small Using interval notation.}

区间表示：$(0 - \delta, 0 + \delta)$

$\quad\Rightarrow\quad (0 - \delta, 0 + \delta)$


\vspace{0.8em}\hrule\vspace{0.8em}

\problem{P3}

\begin{stepbox}
Suppose $a$ is a real number and $\delta$ is a positive real number. How do you describe all real numbers $x$ that are within $\delta$ of $a$, but not equal to $a$:
\end{stepbox}

\solution

\begin{stepbox}
x 与 0 的距离在 δ 以内且不等于 0：$0 < |x - 0| < \delta$
\end{stepbox}

\begin{center}
$\boxed{0 < |x - 0| < \delta}$
\end{center}

\begin{verifybox}\small \verifyok\quad 1/1 项检验通过（答案形式检查） \end{verifybox}

\vspace{0.3em}
\textbf{(a)}\quad \textcolor{gray}{\small Graphically.}

\textcolor{warnamber}{未解：未识别的 ε-δ 子题类型}

\textbf{(b)}\quad \textcolor{gray}{\small Using inequalities.}

用不等式表示：$0 - \delta < x < 0 + \delta, \quad x \neq 0$

$\quad\Rightarrow\quad 0 - \delta < x < 0 + \delta, \quad x \neq 0$

\textbf{(c)}\quad \textcolor{gray}{\small Using absolute value notation.}

$0 < |x - 0| < \delta$

$\quad\Rightarrow\quad 0 < |x - 0| < \delta$

\textbf{(d)}\quad \textcolor{gray}{\small Using interval notation.}

区间表示：$(0 - \delta, 0) \cup (0, 0 + \delta)$

$\quad\Rightarrow\quad (0 - \delta, 0) \cup (0, 0 + \delta)$


\vspace{0.8em}\hrule\vspace{0.8em}

\problem{P4}

\begin{stepbox}
Let $f(x) = x^2$.
\end{stepbox}

\solution

\begin{stepbox}
已知 $f(x) = x^{2}$，中心 $a = 3$
\end{stepbox}

\begin{stepbox}
$f(a) = 9$，故目标值 $L = 9$
\end{stepbox}

\begin{center}
\textbackslash{}text\{见各子题解答\}
\end{center}

\begin{verifybox}\small \verifyok\quad 1/1 项检验通过（答案形式检查） \end{verifybox}

\vspace{0.3em}
\textbf{(a)}\quad \textcolor{gray}{\small Find all the positive numbers \$x\$ such that \$f(x)\$ is within 1 of 9.}

\textcolor{warnamber}{未解：该子题需要解不等式，未覆盖}

\textbf{(b)}\quad \textcolor{gray}{\small Find a real number \$\textbackslash{}delta\$ such that whenever \$x\$ is within \$\textbackslash{}delta\$ of 3, \$f(…}

要使 $|f(x) - 9| < 1$，先作差分解：

$|f(x)-f(a)| = |x-a|\cdot(\text{与 }x\text{ 无关的因子})$

限制 $\delta\le 1$，把后面的因子界为常数 $C$，得

$\delta = \min\{1,\ \varepsilon/C\} = \frac{1}{8}$

验证：取 $\delta = \frac{1}{8}$，当 $0<|x - 3| < \delta$ 时$|f(x) - 9| < 1$ ✓

$\quad\Rightarrow\quad \frac{1}{8}$

\textbf{(c)}\quad \textcolor{gray}{\small Find a real number \$\textbackslash{}delta\$ such that whenever \$x\$ is within \$\textbackslash{}delta\$ of 3, \$f(…}

要使 $|f(x) - 9| < \frac{1}{2}$，先作差分解：

$|f(x)-f(a)| = |x-a|\cdot(\text{与 }x\text{ 无关的因子})$

限制 $\delta\le 1$，把后面的因子界为常数 $C$，得

$\delta = \min\{1,\ \varepsilon/C\} = \frac{1}{16}$

验证：取 $\delta = \frac{1}{16}$，当 $0<|x - 3| < \delta$ 时$|f(x) - 9| < \frac{1}{2}$ ✓

$\quad\Rightarrow\quad \frac{1}{16}$

\textbf{(d)}\quad \textcolor{gray}{\small Is it true that for any positive number \$\textbackslash{}varepsilon\$, there is a positive \$\textbackslash{}de…}

要使 $|f(x) - 9| < \epsilon$，先作差分解：

$|f(x)-f(a)| = |x-a|\cdot(\text{与 }x\text{ 无关的因子})$

限制 $\delta\le 1$，把后面的因子界为常数 $C$，得

$\delta = \min\{1,\ \varepsilon/C\} = \min\left(1, \frac{\epsilon}{8}\right)$

验证：取 $\delta = \min\left(1, \frac{\epsilon}{8}\right)$，当 $0<|x - 3| < \delta$ 时$|f(x) - 9| < \epsilon$ ✓

$\quad\Rightarrow\quad \min\left(1, \frac{\epsilon}{8}\right)$


\vspace{0.8em}\hrule\vspace{0.8em}

\problem{AP1}

\begin{stepbox}
Graph $y = 2x - x^2$. For which points $a$ is the tangent line to the curve at the point $(a, 2a - a^2)$ a horizontal line? How is the value of $2x - x^2$ at these points related to the values of $2x - x^2$ at other points?
\end{stepbox}

\solution

\begin{stepbox}
已知 $y = - x^{2} + 2 x$
\end{stepbox}

\begin{stepbox}
求导：$y' = 2 - 2 x$
\end{stepbox}

\begin{stepbox}
水平切线处 $y' = 0$：解得 $x = ['1']$
\end{stepbox}

\begin{stepbox}
对应点：$(1, 1)$
\end{stepbox}

\begin{stepbox}
$f''(1) = -2 < 0$：极大值点
\end{stepbox}

\begin{center}
$\boxed{(1, 1)}$
\end{center}

\begin{verifybox}\small \verifyok\quad 1/1 项检验通过（答案形式检查） \end{verifybox}

\vspace{0.8em}\hrule\vspace{0.8em}

\problem{AP2}

\begin{stepbox}
The point $(2, 0)$ is on the curve $y = 2x - x^2$.
\end{stepbox}

\solution

\begin{stepbox}
见各子题计算过程。
\end{stepbox}

\begin{center}
\textbackslash{}text\{见各子题\}
\end{center}

\begin{verifybox}\small \verifyok\quad 1/1 项检验通过（答案形式检查） \end{verifybox}

\vspace{0.3em}
\textbf{(a)}\quad \textcolor{gray}{\small Draw the tangent line to the curve at the point \$(2, 0)\$.}

This sub-question asks for a graph; see the matplotlib figure generated for this problem.

$\quad\Rightarrow\quad \text{见附图}$

\textbf{(b)}\quad \textcolor{gray}{\small Try to determine the equation of this tangent line.}

\textcolor{warnamber}{未解：无法解析该子题的极限表达式}


\vspace{0.8em}\hrule\vspace{0.8em}

\vspace{1.5em}
\hrule\vspace{0.8em}
\begin{center}\small
\textbf{求解汇总}\quad 共 8 题，自动解出 8 题（100\%）；验证通过 8 题，存疑 0 题。\\[0.3em]
\textcolor{gray}{求解器分布：conceptual\_template×2、epsilon\_delta\_computation×1、epsilon\_delta\_template×3、sympy\_limit×1、sympy\_tangent×1}
\end{center}
\vfill\begin{center}\footnotesize\textcolor{gray}{由 Meteorain C4C 作业自动求解流水线自动生成}\end{center}
\end{document}
